
\documentclass{article}
\usepackage{colortbl}
\usepackage{makecell}
\usepackage{multirow}
\usepackage{supertabular}

\begin{document}

\newcounter{utterance}

\centering \large Interaction Transcript for game `hot\_air\_balloon', experiment `air\_balloon\_survival\_tr\_negotiation\_easy', episode 0 with Apertus{-}v1.5{-}8B.
\vspace{24pt}

{ \footnotesize  \setcounter{utterance}{1}
\setlength{\tabcolsep}{0pt}
\begin{supertabular}{c@{$\;$}|p{.15\linewidth}@{}p{.15\linewidth}p{.15\linewidth}p{.15\linewidth}p{.15\linewidth}p{.15\linewidth}}
    \# & $\;$A & \multicolumn{4}{c}{Game Master} & $\;\:$B\\
    \hline

    \theutterance \stepcounter{utterance}  
    & & \multicolumn{4}{p{0.6\linewidth}}{
        \cellcolor[rgb]{0.9,0.9,0.9}{
            \makecell[{{p{\linewidth}}}]{
                \texttt{\tiny{[P1$\langle$GM]}}
                \texttt{İşbirliğine dayalı bir müzakere oyununa katılıyorsun.} \\
\\ 
\texttt{Başka bir katılımcıyla birlikte, elde tutulacak tek bir öğe kümesi üzerinde anlaşmanız gerekiyor. Her birinizin, her bir öğenin sizin için ne kadar önemli olduğuna dair kendi bakışı (önem değeri) vardır. Diğer katılımcının öğelere nasıl değer verdiğini bilmiyorsunuz. Ayrıca, her bir öğenin gerektirdiği çaba değeri de size verilmiştir.} \\
\texttt{Yalnızca seçilen öğelerin toplam çabası paylaşılan bir sınırı aşmıyorsa bir küme üzerinde anlaşabilirsiniz:} \\
\\ 
\texttt{LİMİT = 3358} \\
\\ 
\texttt{İşte öğelerin tek tek çaba değerleri:} \\
\\ 
\texttt{Öğe çabaları = \{"C08": 171, "C55": 933, "A43": 552, "A67": 219, "B42": 275, "A62": 778, "B19": 341, "C43": 615, "C73": 519, "C10": 862, "B36": 262, "A18": 377, "C61": 347, "C67": 349, "C42": 117\}} \\
\\ 
\texttt{İşte her bir öğenin önemine dair kişisel bakışın:} \\
\\ 
\texttt{Öğe önem değerleri = \{"C08": 138, "C55": 583, "A43": 868, "A67": 822, "B42": 783, "A62": 65, "B19": 262, "C43": 121, "C73": 508, "C10": 780, "B36": 461, "A18": 484, "C61": 668, "C67": 389, "C42": 808\}} \\
\\ 
\texttt{Amaç:} \\
\\ 
\texttt{Amacın, LİMİT içinde kalarak sana mümkün olan en fazla faydayı sağlayan (yani SENİN için toplam önemi maksimize eden) ortak bir öğe kümesi üzerinde anlaşmaktır. Her mesajda mutlaka bir TEKLİF yapmak zorunda değilsin {-} sadece müzakere de edebilirsin. Tüm taktikler serbest!} \\
\\ 
\texttt{Etkileşim Protokolü:} \\
\\ 
\texttt{Bir mesajda yalnızca aşağıdaki yapılandırılmış biçimleri kullanabilirsin:} \\
\\ 
\texttt{TEKLİF: \{'A', 'B', 'C', …\}} \\
\texttt{Belirtilen öğelerin tutulmasını teklif eder.} \\
\texttt{REDDET: \{'A', 'B', 'C', …\}} \\
\texttt{Rakibin teklifini açıkça reddeder.} \\
\texttt{SAVUNMA: \{'...'\}} \\
\texttt{Son teklifinizi savunur veya diğer oyuncunun teklifine karşı argüman sunar.} \\
\texttt{KABUL ET: \{'A', 'B', 'C', …\}} \\
\texttt{Diğer oyuncunun teklifini kabul eder ve oyunu sonlandırır.} \\
\\ 
\\ 
\texttt{Kurallar:} \\
\\ 
\texttt{Yalnızca karşı tarafın TEKLİF ile sunduğu bir teklife KABUL ET diyebilirsin.} \\
\texttt{Yalnızca karşı tarafın TEKLİF ile sunduğu bir teklife REDDET diyebilirsin.} \\
\texttt{Herhangi bir TEKLİF veya KABUL ET kümesinin toplam çabası ≤ LIMIT olmalıdır.} \\
\texttt{Gizli önem değerlerini AÇIKLAMA.} \\
\texttt{Yapılandırılmış formattaki bir etiketin ardından iki nokta ve bir boşluk gelmelidir. İçerik, 0 veya daha fazla string içeren bir Python kümesi olmalıdır.} \\
\texttt{Yani format şu şekilde olmalıdır: "ETİKET: \{...\}"} \\
\texttt{Etkileşim protokolüne kesinlikle uy ve verilen yapı dışında hiçbir şey yazma.} \\
\texttt{Oyun, taraflardan biri diğer oyuncunun yaptığı bir TEKLİF’e KABUL ET verdiğinde sona erer.} \\
\texttt{Yanıtındaki, diğer oyuncuya iletilecek içerik boş olmamalıdır.} \\
\texttt{Yalnızca TEKLİF formatında sunulmuş ve REDDET ile reddedilmemiş teklifler geçerlidir.} \\
\\ 
            }
        }
    }
    & \\ \\

    \theutterance \stepcounter{utterance}  
    & \multicolumn{4}{p{0.6\linewidth}}{
        \cellcolor[rgb]{0.9,0.9,0.9}{
            \makecell[{{p{\linewidth}}}]{
                \texttt{\tiny{[P1$\rangle$GM]}}
                \texttt{TEKLİF: \{'C08', 'C42'\}} \\
            }
        }
    }
    & & \\ \\

    \theutterance \stepcounter{utterance}  
    & & \multicolumn{4}{p{0.6\linewidth}}{
        \cellcolor[rgb]{0.9,0.9,0.9}{
            \makecell[{{p{\linewidth}}}]{
                \texttt{\tiny{[GM$\rangle$P2]}}
                \texttt{İşbirliğine dayalı bir müzakere oyununa katılıyorsun.} \\
\\ 
\texttt{Başka bir katılımcıyla birlikte, elde tutulacak tek bir öğe kümesi üzerinde anlaşmanız gerekiyor. Her birinizin, her bir öğenin sizin için ne kadar önemli olduğuna dair kendi bakışı (önem değeri) vardır. Diğer katılımcının öğelere nasıl değer verdiğini bilmiyorsunuz. Ayrıca, her bir öğenin gerektirdiği çaba değeri de size verilmiştir.} \\
\texttt{Yalnızca seçilen öğelerin toplam çabası paylaşılan bir sınırı aşmıyorsa bir küme üzerinde anlaşabilirsiniz:} \\
\\ 
\texttt{LİMİT = 3358} \\
\\ 
\texttt{İşte öğelerin tek tek çaba değerleri:} \\
\\ 
\texttt{Öğe çabaları = \{"C08": 171, "C55": 933, "A43": 552, "A67": 219, "B42": 275, "A62": 778, "B19": 341, "C43": 615, "C73": 519, "C10": 862, "B36": 262, "A18": 377, "C61": 347, "C67": 349, "C42": 117\}} \\
\\ 
\texttt{İşte her bir öğenin önemine dair kişisel bakışın:} \\
\\ 
\texttt{Öğe önem değerleri = \{"C08": 138, "C55": 583, "A43": 868, "A67": 822, "B42": 783, "A62": 65, "B19": 262, "C43": 121, "C73": 508, "C10": 780, "B36": 461, "A18": 484, "C61": 668, "C67": 389, "C42": 808\}} \\
\\ 
\texttt{Amaç:} \\
\\ 
\texttt{Amacın, LİMİT içinde kalarak sana mümkün olan en fazla faydayı sağlayan (yani SENİN için toplam önemi maksimize eden) ortak bir öğe kümesi üzerinde anlaşmaktır. Her mesajda mutlaka bir TEKLİF yapmak zorunda değilsin {-} sadece müzakere de edebilirsin. Tüm taktikler serbest!} \\
\\ 
\texttt{Etkileşim Protokolü:} \\
\\ 
\texttt{Bir mesajda yalnızca aşağıdaki yapılandırılmış biçimleri kullanabilirsin:} \\
\\ 
\texttt{TEKLİF: \{'A', 'B', 'C', …\}} \\
\texttt{Belirtilen öğelerin tutulmasını teklif eder.} \\
\texttt{REDDET: \{'A', 'B', 'C', …\}} \\
\texttt{Rakibin teklifini açıkça reddeder.} \\
\texttt{SAVUNMA: \{'...'\}} \\
\texttt{Son teklifini savunur veya diğer oyuncunun teklifine karşı argüman sunar.} \\
\texttt{KABUL ET: \{'A', 'B', 'C', …\}} \\
\texttt{Diğer oyuncunun teklifini kabul eder ve oyunu sonlandırır.} \\
\\ 
\\ 
\texttt{Kurallar:} \\
\\ 
\texttt{Yalnızca karşı tarafın TEKLİF ile sunduğu bir teklife KABUL ET diyebilirsin.} \\
\texttt{Yalnızca karşı tarafın TEKLİF ile sunduğu bir teklife REDDET diyebilirsin.} \\
\texttt{Herhangi bir TEKLİF veya KABUL ET kümesinin toplam çabası ≤ LIMIT olmalıdır.} \\
\texttt{Gizli önem değerlerini AÇIKLAMA.} \\
\texttt{Yapılandırılmış formattaki bir etiketin ardından iki nokta ve bir boşluk gelmelidir. İçerik, 0 veya daha fazla string içeren bir Python kümesi olmalıdır.} \\
\texttt{Yani format şu şekilde olmalıdır: "ETİKET: \{...\}"} \\
\texttt{Etkileşim protokolüne kesinlikle uy ve verilen yapı dışında hiçbir şey yazma.} \\
\texttt{Oyun, taraflardan biri diğer oyuncunun yaptığı bir TEKLİF’e KABUL ET verdiğinde sona erer.} \\
\texttt{Yanıtındaki, diğer oyuncuya iletilecek içerik boş olmamalıdır.} \\
\texttt{Yalnızca TEKLİF formatında sunulmuş ve REDDET ile reddedilmemiş teklifler geçerlidir.} \\
\\ 
\\ 
\\ 
\texttt{Şimdi diğer oyuncunun ilk mesajını alacaksın.} \\
\\ 
\texttt{TEKLİF: \{'C08', 'C42'\}} \\
            }
        }
    }
    & \\ \\

    \theutterance \stepcounter{utterance}  
    & & & \multicolumn{4}{p{0.6\linewidth}}{
        \cellcolor[rgb]{0.9,0.9,0.9}{
            \makecell[{{p{\linewidth}}}]{
                \texttt{\tiny{[GM$\langle$P2]}}
                \texttt{TEKLİF: \{'C08', 'C42'\}} \\
            }
        }
    }
     \\ \\

    \theutterance \stepcounter{utterance}  
    & & \multicolumn{4}{p{0.6\linewidth}}{
        \cellcolor[rgb]{0.9,0.9,0.9}{
            \makecell[{{p{\linewidth}}}]{
                \texttt{\tiny{[P1$\langle$GM]}}
                \texttt{TEKLİF: \{'C08', 'C42'\}} \\
            }
        }
    }
    & \\ \\

    \theutterance \stepcounter{utterance}  
    & \multicolumn{4}{p{0.6\linewidth}}{
        \cellcolor[rgb]{0.9,0.9,0.9}{
            \makecell[{{p{\linewidth}}}]{
                \texttt{\tiny{[P1$\rangle$GM]}}
                \texttt{SAVUNMA: Önerilen öğelerin toplam çabası 171 + 117 = 288'dir ki limit olan 3358'den çok uzakta kalınmaktadır. Bu öğelerin toplam önemi ise 138 + 808 = 946'dır. Öncelikle düşük çaba gerektiren ve önemli öğelerden oluşması açısından mantıklı bir başlangıç teklifi olarak değerlendirilebilir. Ancak kümenin büyüklenme potansiyeli ve diğer öğelerin önemi göz önüne alındığında, müzakere devam edilmelidir.} \\
            }
        }
    }
    & & \\ \\

    \theutterance \stepcounter{utterance}  
    & & \multicolumn{4}{p{0.6\linewidth}}{
        \cellcolor[rgb]{0.9,0.9,0.9}{
            \makecell[{{p{\linewidth}}}]{
                \texttt{\tiny{[GM$\rangle$P2]}}
                \texttt{SAVUNMA: Önerilen öğelerin toplam çabası 171 + 117 = 288'dir ki limit olan 3358'den çok uzakta kalınmaktadır. Bu öğelerin toplam önemi ise 138 + 808 = 946'dır. Öncelikle düşük çaba gerektiren ve önemli öğelerden oluşması açısından mantıklı bir başlangıç teklifi olarak değerlendirilebilir. Ancak kümenin büyüklenme potansiyeli ve diğer öğelerin önemi göz önüne alındığında, müzakere devam edilmelidir.} \\
            }
        }
    }
    & \\ \\

    \theutterance \stepcounter{utterance}  
    & & & \multicolumn{4}{p{0.6\linewidth}}{
        \cellcolor[rgb]{0.9,0.9,0.9}{
            \makecell[{{p{\linewidth}}}]{
                \texttt{\tiny{[GM$\langle$P2]}}
                \texttt{SAVUNMA: Önerilen öğelerin toplam çabası 288, toplam önemi 946'dır. Bu öğeler düşük çaba ve belirli bir önem değerine sahiptir. Ancak kümenin büyüklenmesi ve diğer öğelerin önemi göz önüne alındığında, diğer öğelerin değerlendirilmesi gerekmektedir. Örneğin, yüksek öneme sahip ancak yüksek çaba gerektiren öğeler de değerlendirilebilir. Müzakere devam edebilir.} \\
            }
        }
    }
     \\ \\

    \theutterance \stepcounter{utterance}  
    & & \multicolumn{4}{p{0.6\linewidth}}{
        \cellcolor[rgb]{0.9,0.9,0.9}{
            \makecell[{{p{\linewidth}}}]{
                \texttt{\tiny{[P1$\langle$GM]}}
                \texttt{SAVUNMA: Önerilen öğelerin toplam çabası 288, toplam önemi 946'dır. Bu öğeler düşük çaba ve belirli bir önem değerine sahiptir. Ancak kümenin büyüklenmesi ve diğer öğelerin önemi göz önüne alındığında, diğer öğelerin değerlendirilmesi gerekmektedir. Örneğin, yüksek öneme sahip ancak yüksek çaba gerektiren öğeler de değerlendirilebilir. Müzakere devam edebilir.} \\
            }
        }
    }
    & \\ \\

    \theutterance \stepcounter{utterance}  
    & \multicolumn{4}{p{0.6\linewidth}}{
        \cellcolor[rgb]{0.9,0.9,0.9}{
            \makecell[{{p{\linewidth}}}]{
                \texttt{\tiny{[P1$\rangle$GM]}}
                \texttt{TEKLİF: \{'C10', 'A43'\}} \\
            }
        }
    }
    & & \\ \\

    \theutterance \stepcounter{utterance}  
    & & \multicolumn{4}{p{0.6\linewidth}}{
        \cellcolor[rgb]{0.9,0.9,0.9}{
            \makecell[{{p{\linewidth}}}]{
                \texttt{\tiny{[GM$\rangle$P2]}}
                \texttt{TEKLİF: \{'C10', 'A43'\}} \\
            }
        }
    }
    & \\ \\

    \theutterance \stepcounter{utterance}  
    & & & \multicolumn{4}{p{0.6\linewidth}}{
        \cellcolor[rgb]{0.9,0.9,0.9}{
            \makecell[{{p{\linewidth}}}]{
                \texttt{\tiny{[GM$\langle$P2]}}
                \texttt{SAVUNMA: Önerilen öğelerin toplam çabası 862 + 552 = 1414'tür. Toplam önemi ise 780 + 868 = 1648'dir. Bu öğeler, yüksek önem değerlerine sahip olup orta düzeyde çaba gerektirmektedir. Ancak, kümenin toplam çabasının limit olan 3358'den daha az kalması beklenirken, yüksek öneme sahip öğelerin seçilmesi stratejik olarak doğru bir hamle sayılabilir. Ancak, diğer olası öğeler ve onların önemine değinerek, müzakere devam edilmelidir.} \\
            }
        }
    }
     \\ \\

    \theutterance \stepcounter{utterance}  
    & & \multicolumn{4}{p{0.6\linewidth}}{
        \cellcolor[rgb]{0.9,0.9,0.9}{
            \makecell[{{p{\linewidth}}}]{
                \texttt{\tiny{[P1$\langle$GM]}}
                \texttt{SAVUNMA: Önerilen öğelerin toplam çabası 862 + 552 = 1414'tür. Toplam önemi ise 780 + 868 = 1648'dir. Bu öğeler, yüksek önem değerlerine sahip olup orta düzeyde çaba gerektirmektedir. Ancak, kümenin toplam çabasının limit olan 3358'den daha az kalması beklenirken, yüksek öneme sahip öğelerin seçilmesi stratejik olarak doğru bir hamle sayılabilir. Ancak, diğer olası öğeler ve onların önemine değinerek, müzakere devam edilmelidir.} \\
            }
        }
    }
    & \\ \\

    \theutterance \stepcounter{utterance}  
    & \multicolumn{4}{p{0.6\linewidth}}{
        \cellcolor[rgb]{0.9,0.9,0.9}{
            \makecell[{{p{\linewidth}}}]{
                \texttt{\tiny{[P1$\rangle$GM]}}
                \texttt{TEKLİF: \{'C55', 'B42'\}} \\
            }
        }
    }
    & & \\ \\

    \theutterance \stepcounter{utterance}  
    & & \multicolumn{4}{p{0.6\linewidth}}{
        \cellcolor[rgb]{0.9,0.9,0.9}{
            \makecell[{{p{\linewidth}}}]{
                \texttt{\tiny{[GM$\rangle$P2]}}
                \texttt{TEKLİF: \{'C55', 'B42'\}} \\
            }
        }
    }
    & \\ \\

    \theutterance \stepcounter{utterance}  
    & & & \multicolumn{4}{p{0.6\linewidth}}{
        \cellcolor[rgb]{0.9,0.9,0.9}{
            \makecell[{{p{\linewidth}}}]{
                \texttt{\tiny{[GM$\langle$P2]}}
                \texttt{SAVUNMA: Önerilen öğelerin toplam çabası 933 + 275 = 1208'dir. Toplam önemi ise 583 + 783 = 1366'dır. Bu öğeler, orta düzeyde çaba gerektiren ve belirli önem değerine sahip öğelerdir. Ancak, kümenin toplam çabasının limitten uzakta kalması gerekmektedir ve daha fazla öğe eklenebilir. Örneğin, önemi yüksek ancak çaba açısından daha düşük olan öğelerin değerlendirilmesi ile kümenin genişletilmesi mümkündür. Müzakere devam etmeli.} \\
            }
        }
    }
     \\ \\

    \theutterance \stepcounter{utterance}  
    & & \multicolumn{4}{p{0.6\linewidth}}{
        \cellcolor[rgb]{0.9,0.9,0.9}{
            \makecell[{{p{\linewidth}}}]{
                \texttt{\tiny{[P1$\langle$GM]}}
                \texttt{SAVUNMA: Önerilen öğelerin toplam çabası 933 + 275 = 1208'dir. Toplam önemi ise 583 + 783 = 1366'dır. Bu öğeler, orta düzeyde çaba gerektiren ve belirli önem değerine sahip öğelerdir. Ancak, kümenin toplam çabasının limitten uzakta kalması gerekmektedir ve daha fazla öğe eklenebilir. Örneğin, önemi yüksek ancak çaba açısından daha düşük olan öğelerin değerlendirilmesi ile kümenin genişletilmesi mümkündür. Müzakere devam etmeli.} \\
            }
        }
    }
    & \\ \\

    \theutterance \stepcounter{utterance}  
    & \multicolumn{4}{p{0.6\linewidth}}{
        \cellcolor[rgb]{0.9,0.9,0.9}{
            \makecell[{{p{\linewidth}}}]{
                \texttt{\tiny{[P1$\rangle$GM]}}
                \texttt{TEKLİF: \{'A67', 'B19'\}} \\
            }
        }
    }
    & & \\ \\

    \theutterance \stepcounter{utterance}  
    & & \multicolumn{4}{p{0.6\linewidth}}{
        \cellcolor[rgb]{0.9,0.9,0.9}{
            \makecell[{{p{\linewidth}}}]{
                \texttt{\tiny{[GM$\rangle$P2]}}
                \texttt{TEKLİF: \{'A67', 'B19'\}} \\
            }
        }
    }
    & \\ \\

    \theutterance \stepcounter{utterance}  
    & & & \multicolumn{4}{p{0.6\linewidth}}{
        \cellcolor[rgb]{0.9,0.9,0.9}{
            \makecell[{{p{\linewidth}}}]{
                \texttt{\tiny{[GM$\langle$P2]}}
                \texttt{SAVUNMA: Önerilen öğelerin toplam çabası 219 + 341 = 560'dır. Toplam önemi ise 822 + 262 = 1084'dir. Bu öğeler düşük çaba gerektiren ve belirli önem değerleri taşıyan öğelerdir. Ancak, kümenin toplam çabasının limitten uzakta kalması için büyüklenme potansiyeli olan daha yüksek öneme sahip öğelerin değerlendirilmesi gerekmektedir. Müzakere devam edilmelidir.} \\
            }
        }
    }
     \\ \\

    \theutterance \stepcounter{utterance}  
    & & \multicolumn{4}{p{0.6\linewidth}}{
        \cellcolor[rgb]{0.9,0.9,0.9}{
            \makecell[{{p{\linewidth}}}]{
                \texttt{\tiny{[P1$\langle$GM]}}
                \texttt{SAVUNMA: Önerilen öğelerin toplam çabası 219 + 341 = 560'dır. Toplam önemi ise 822 + 262 = 1084'dir. Bu öğeler düşük çaba gerektiren ve belirli önem değerleri taşıyan öğelerdir. Ancak, kümenin toplam çabasının limitten uzakta kalması için büyüklenme potansiyeli olan daha yüksek öneme sahip öğelerin değerlendirilmesi gerekmektedir. Müzakere devam edilmelidir.} \\
            }
        }
    }
    & \\ \\

    \theutterance \stepcounter{utterance}  
    & \multicolumn{4}{p{0.6\linewidth}}{
        \cellcolor[rgb]{0.9,0.9,0.9}{
            \makecell[{{p{\linewidth}}}]{
                \texttt{\tiny{[P1$\rangle$GM]}}
                \texttt{TEKLİF: \{'A62', 'C61'\}} \\
            }
        }
    }
    & & \\ \\

    \theutterance \stepcounter{utterance}  
    & & \multicolumn{4}{p{0.6\linewidth}}{
        \cellcolor[rgb]{0.9,0.9,0.9}{
            \makecell[{{p{\linewidth}}}]{
                \texttt{\tiny{[GM$\rangle$P2]}}
                \texttt{TEKLİF: \{'A62', 'C61'\}} \\
            }
        }
    }
    & \\ \\

    \theutterance \stepcounter{utterance}  
    & & & \multicolumn{4}{p{0.6\linewidth}}{
        \cellcolor[rgb]{0.9,0.9,0.9}{
            \makecell[{{p{\linewidth}}}]{
                \texttt{\tiny{[GM$\langle$P2]}}
                \texttt{SAVUNMA: Önerilen öğelerin toplam çabası 778 + 347 = 1125'dir. Toplam önemi ise 65 + 668 = 733'dir. Bu öğelerin toplam önemi düşük bir değere sahipken, toplam çaba değeri orta seviyededir. Kümenin stratejik olarak daha yüksek öneme sahip öğelerle zenginleştirilmesi yönünde hareket edilmelidir. Alternatif öğelerin önemi ve çabası göz önüne alındığında, müzakere devam etmelidir.} \\
            }
        }
    }
     \\ \\

    \theutterance \stepcounter{utterance}  
    & & \multicolumn{4}{p{0.6\linewidth}}{
        \cellcolor[rgb]{0.9,0.9,0.9}{
            \makecell[{{p{\linewidth}}}]{
                \texttt{\tiny{[P1$\langle$GM]}}
                \texttt{SAVUNMA: Önerilen öğelerin toplam çabası 778 + 347 = 1125'dir. Toplam önemi ise 65 + 668 = 733'dir. Bu öğelerin toplam önemi düşük bir değere sahipken, toplam çaba değeri orta seviyededir. Kümenin stratejik olarak daha yüksek öneme sahip öğelerle zenginleştirilmesi yönünde hareket edilmelidir. Alternatif öğelerin önemi ve çabası göz önüne alındığında, müzakere devam etmelidir.} \\
            }
        }
    }
    & \\ \\

    \theutterance \stepcounter{utterance}  
    & \multicolumn{4}{p{0.6\linewidth}}{
        \cellcolor[rgb]{0.9,0.9,0.9}{
            \makecell[{{p{\linewidth}}}]{
                \texttt{\tiny{[P1$\rangle$GM]}}
                \texttt{TEKLİF: \{'C67', 'B36'\}} \\
            }
        }
    }
    & & \\ \\

    \theutterance \stepcounter{utterance}  
    & & \multicolumn{4}{p{0.6\linewidth}}{
        \cellcolor[rgb]{0.9,0.9,0.9}{
            \makecell[{{p{\linewidth}}}]{
                \texttt{\tiny{[GM$\rangle$P2]}}
                \texttt{TEKLİF: \{'C67', 'B36'\}} \\
            }
        }
    }
    & \\ \\

    \theutterance \stepcounter{utterance}  
    & & & \multicolumn{4}{p{0.6\linewidth}}{
        \cellcolor[rgb]{0.9,0.9,0.9}{
            \makecell[{{p{\linewidth}}}]{
                \texttt{\tiny{[GM$\langle$P2]}}
                \texttt{SAVUNMA: Önerilen öğenin toplam çabası 349 + 262 = 611'dir. Toplam önemi ise 389 + 461 = 850'dir. Bu öğelerin toplam önemi orta seviyededir ve toplam çaba değeri orta seviyededir. Ancak kümenin total önemi maksimize edilmesi gerekmektedir ve daha yüksek öneme sahip öğelerin değerlendirilmesi stratejik olarak önemlidir. Müzakere devam edilmelidir.} \\
            }
        }
    }
     \\ \\

    \theutterance \stepcounter{utterance}  
    & & \multicolumn{4}{p{0.6\linewidth}}{
        \cellcolor[rgb]{0.9,0.9,0.9}{
            \makecell[{{p{\linewidth}}}]{
                \texttt{\tiny{[P1$\langle$GM]}}
                \texttt{SAVUNMA: Önerilen öğenin toplam çabası 349 + 262 = 611'dir. Toplam önemi ise 389 + 461 = 850'dir. Bu öğelerin toplam önemi orta seviyededir ve toplam çaba değeri orta seviyededir. Ancak kümenin total önemi maksimize edilmesi gerekmektedir ve daha yüksek öneme sahip öğelerin değerlendirilmesi stratejik olarak önemlidir. Müzakere devam edilmelidir.} \\
            }
        }
    }
    & \\ \\

    \theutterance \stepcounter{utterance}  
    & \multicolumn{4}{p{0.6\linewidth}}{
        \cellcolor[rgb]{0.9,0.9,0.9}{
            \makecell[{{p{\linewidth}}}]{
                \texttt{\tiny{[P1$\rangle$GM]}}
                \texttt{TEKLİF: \{'A18', 'C73'\}} \\
            }
        }
    }
    & & \\ \\

    \theutterance \stepcounter{utterance}  
    & & \multicolumn{4}{p{0.6\linewidth}}{
        \cellcolor[rgb]{0.9,0.9,0.9}{
            \makecell[{{p{\linewidth}}}]{
                \texttt{\tiny{[GM$\rangle$P2]}}
                \texttt{TEKLİF: \{'A18', 'C73'\}} \\
            }
        }
    }
    & \\ \\

    \theutterance \stepcounter{utterance}  
    & & & \multicolumn{4}{p{0.6\linewidth}}{
        \cellcolor[rgb]{0.9,0.9,0.9}{
            \makecell[{{p{\linewidth}}}]{
                \texttt{\tiny{[GM$\langle$P2]}}
                \texttt{SAVUNMA: Önerilen öğelerin toplam çabası 377 + 508 = 885'dir. Toplam önemi ise 484 + 508 = 992'dir. Bu öğelerin toplam önemi yüksektir ve toplam çaba değeri orta seviyededir. Kümenin büyüklenmesi için yüksek öneme sahip olan ancak daha düşük çaba gerektiren öğelerin değerlendirilmesi gerekmektedir. Müzakere devam edilmelidir.} \\
            }
        }
    }
     \\ \\

    \theutterance \stepcounter{utterance}  
    & & \multicolumn{4}{p{0.6\linewidth}}{
        \cellcolor[rgb]{0.9,0.9,0.9}{
            \makecell[{{p{\linewidth}}}]{
                \texttt{\tiny{[P1$\langle$GM]}}
                \texttt{SAVUNMA: Önerilen öğelerin toplam çabası 377 + 508 = 885'dir. Toplam önemi ise 484 + 508 = 992'dir. Bu öğelerin toplam önemi yüksektir ve toplam çaba değeri orta seviyededir. Kümenin büyüklenmesi için yüksek öneme sahip olan ancak daha düşük çaba gerektiren öğelerin değerlendirilmesi gerekmektedir. Müzakere devam edilmelidir.} \\
            }
        }
    }
    & \\ \\

    \theutterance \stepcounter{utterance}  
    & \multicolumn{4}{p{0.6\linewidth}}{
        \cellcolor[rgb]{0.9,0.9,0.9}{
            \makecell[{{p{\linewidth}}}]{
                \texttt{\tiny{[P1$\rangle$GM]}}
                \texttt{TEKLİF: \{'C43'\}} \\
            }
        }
    }
    & & \\ \\

    \theutterance \stepcounter{utterance}  
    & & \multicolumn{4}{p{0.6\linewidth}}{
        \cellcolor[rgb]{0.9,0.9,0.9}{
            \makecell[{{p{\linewidth}}}]{
                \texttt{\tiny{[GM$\rangle$P2]}}
                \texttt{TEKLİF: \{'C43'\}} \\
            }
        }
    }
    & \\ \\

    \theutterance \stepcounter{utterance}  
    & & & \multicolumn{4}{p{0.6\linewidth}}{
        \cellcolor[rgb]{0.9,0.9,0.9}{
            \makecell[{{p{\linewidth}}}]{
                \texttt{\tiny{[GM$\langle$P2]}}
                \texttt{SAVUNMA: Önerilen öğenin toplam çabası 615'dir. Toplam önemi ise 121'dir. Bu öğe yalnızca düşük çaba ve öneme sahiptir. Kümenin stratejik olarak daha yüksek öneme sahip ve çaba açısından daha verimli öğelerle zenginleştirilmesi gerekmektedir. Müzakere devam etmelidir.} \\
            }
        }
    }
     \\ \\

    \theutterance \stepcounter{utterance}  
    & & \multicolumn{4}{p{0.6\linewidth}}{
        \cellcolor[rgb]{0.9,0.9,0.9}{
            \makecell[{{p{\linewidth}}}]{
                \texttt{\tiny{[P1$\langle$GM]}}
                \texttt{SAVUNMA: Önerilen öğenin toplam çabası 615'dir. Toplam önemi ise 121'dir. Bu öğe yalnızca düşük çaba ve öneme sahiptir. Kümenin stratejik olarak daha yüksek öneme sahip ve çaba açısından daha verimli öğelerle zenginleştirilmesi gerekmektedir. Müzakere devam etmelidir.} \\
            }
        }
    }
    & \\ \\

    \theutterance \stepcounter{utterance}  
    & \multicolumn{4}{p{0.6\linewidth}}{
        \cellcolor[rgb]{0.9,0.9,0.9}{
            \makecell[{{p{\linewidth}}}]{
                \texttt{\tiny{[P1$\rangle$GM]}}
                \texttt{TEKLİF: \{'B42', 'C10'\}} \\
            }
        }
    }
    & & \\ \\

    \theutterance \stepcounter{utterance}  
    & & \multicolumn{4}{p{0.6\linewidth}}{
        \cellcolor[rgb]{0.9,0.9,0.9}{
            \makecell[{{p{\linewidth}}}]{
                \texttt{\tiny{[GM$\rangle$P2]}}
                \texttt{TEKLİF: \{'B42', 'C10'\}} \\
            }
        }
    }
    & \\ \\

    \theutterance \stepcounter{utterance}  
    & & & \multicolumn{4}{p{0.6\linewidth}}{
        \cellcolor[rgb]{0.9,0.9,0.9}{
            \makecell[{{p{\linewidth}}}]{
                \texttt{\tiny{[GM$\langle$P2]}}
                \texttt{SAVUNMA: Önerilen öğelerin toplam çabası 275 + 862 = 1137'dir. Toplam önemi ise 783 + 780 = 1563'dir. Bu öğelerin yüksek önemi ve orta düzeyde çaba gerektirmesi, kümenin büyüklenmesi için önemli bir fırsat sunmaktadır. Ancak limit değerine uygun olarak daha fazla öğe eklenme potansiyeli değerlendirilmelidir. Müzakere devam edilmelidir.} \\
            }
        }
    }
     \\ \\

    \theutterance \stepcounter{utterance}  
    & & \multicolumn{4}{p{0.6\linewidth}}{
        \cellcolor[rgb]{0.9,0.9,0.9}{
            \makecell[{{p{\linewidth}}}]{
                \texttt{\tiny{[P1$\langle$GM]}}
                \texttt{SAVUNMA: Önerilen öğelerin toplam çabası 275 + 862 = 1137'dir. Toplam önemi ise 783 + 780 = 1563'dir. Bu öğelerin yüksek önemi ve orta düzeyde çaba gerektirmesi, kümenin büyüklenmesi için önemli bir fırsat sunmaktadır. Ancak limit değerine uygun olarak daha fazla öğe eklenme potansiyeli değerlendirilmelidir. Müzakere devam edilmelidir.} \\
            }
        }
    }
    & \\ \\

    \theutterance \stepcounter{utterance}  
    & \multicolumn{4}{p{0.6\linewidth}}{
        \cellcolor[rgb]{0.9,0.9,0.9}{
            \makecell[{{p{\linewidth}}}]{
                \texttt{\tiny{[P1$\rangle$GM]}}
                \texttt{TEKLİF: \{'C61', 'A43'\}} \\
            }
        }
    }
    & & \\ \\

    \theutterance \stepcounter{utterance}  
    & & & \multicolumn{2}{p{0.3\linewidth}}{
        \cellcolor[rgb]{0.9,0.9,0.9}{
            \makecell[{{p{\linewidth}}}]{
                \texttt{\tiny{[GM$|$GM]}}
                \texttt{Maximum rounds exceeded: 10 >= 10} \\
            }
        }
    }
    & & \\ \\

    \theutterance \stepcounter{utterance}  
    & & & \multicolumn{2}{p{0.3\linewidth}}{
        \cellcolor[rgb]{0.9,0.9,0.9}{
            \makecell[{{p{\linewidth}}}]{
                \texttt{\tiny{[GM$|$GM]}}
                \texttt{end game} \\
            }
        }
    }
    & & \\ \\

\end{supertabular}
}

\end{document}
